\subsection{Common shifts and independent classes}
\label{sec:translations}
Let $G$ now be the binary Goldbach ring on a bounded-below support
$\A$, and use $K=K(G)$ for its coefficient Koszul complex.
The upper quadrilateral of Example~\ref{ex:weight33} is one instance
of a construction that works in every monomial degree.
Suppose $\sigma=\{a_1<\cdots<a_d\}$ and $\sigma+h$ both lie in the support
$\A$, for a positive integer $h$. The two lists need not be consecutive
terms of an arithmetic progression. For the binary ring consider
\begin{equation}\label{eq:translated-cycle}
 z_{\sigma,h}=\sum_{j=1}^d(-1)^{j-1}x_{a_j+h}
       \varepsilon_{\sigma\setminus\{a_j\}}.
\end{equation}

\Needspace{8\baselineskip}
\begin{theorem}[Independent translation classes]\label{thm:translation}
Fix $h>0$ and $d\ge2$. As $\sigma$ varies over the $d$-element subsets of
$\A\cap(\A-h)$, the chains $z_{\sigma,h}$ give linearly independent classes
over every field. For each $\sigma$ the class has homological degree $d-1$
and bidegree $(n,d)$, where $n=h+\sum_{a\in \sigma}a$. In particular,
\[
 \beta_{d-1,n,d}\ \ge\
 \#\{\sigma\subseteq\A\cap(\A-h): |\sigma|=d,\ h+\textstyle\sum_{a\in \sigma}a=n\}.
\]
These classes remain independent under enlargement of the support.
\end{theorem}
\begin{proof}
For $j<r$, the terms in the differential omitting both exterior
indices have opposite signs and coefficients $x_{a_j+h}x_{a_r}$
and $x_{a_j}x_{a_r+h}$. Their pair sums agree, so their
difference vanishes in $G$.

Suppose that a nonzero finite combination of the displayed cycles
were a boundary. Let $a$ be the largest label in any $\sigma$ with nonzero
coefficient, and choose such a $\sigma$ containing $a$. The coordinate
$x_{a+h}\varepsilon_{\sigma\setminus\{a\}}$ determines $\sigma$ for the fixed
$h$ and therefore has nonzero coefficient. Any preimage has monomial degree
and homological degree $d$, so its coefficients are scalars. The only
column giving the chosen coordinate is $\mathrm{d}(\varepsilon_\tau)$, where $\tau=(\sigma\setminus\{a\})\cup\{a+h\}$.
Since $d\ge2$, this column has another term with $a+h$ in its
exterior part. No selected cycle has such a term, and no other column
can cancel it: its coefficient label and its exterior labels together
determine $\tau$. The variables of $G$ are independent in monomial degree one,
so this is a contradiction. The proof uses only coefficients $\pm1$
and applies in every larger support.
\end{proof}

\subsection{Consequences for projective dimension}
The independence theorem turns repeated additive differences into
lower bounds for the size and length of a resolution.

\begin{corollary}\label{cor:density}
If the positive differences in $\A$ have unbounded multiplicities,
the binary Goldbach ring on $\A$ has infinite projective dimension over
$R$. In fact each positive homological degree of the linear
strand then has infinitely many generators. Both statements apply to
the positive odd primes.
\end{corollary}
\begin{proof}
An overlap of size $r$ at difference $h$ gives $\binom{r}{i+1}$
independent classes in homological degree $i$. For each fixed $i$ this
becomes arbitrarily large as $r$ increases. On the odd primes, put
$\A_m=\{p\le m:p\text{ odd prime}\}$. Its $\binom{|\A_m|}{2}$
unordered pairs have fewer than $m$ distinct positive differences. Some
difference therefore occurs at least $\binom{|\A_m|}{2}/m$ times,
which tends to infinity by the prime number theorem. Alternatively,
the prime progressions of Green and Tao \cite{GreenTao08} give
arbitrarily large overlaps by shifting each progression one step.
\end{proof}

Appendix~\ref{app:scroll-shift} gives a stronger bound for each fixed
shift using a scroll resolution. In Section~\ref{sec:betti-euler-shifts}
we fix the Goldbach degree instead and select independent classes
whose positive shifts may vary.
